The previous section leaves every cohort in the floor-zero regime, so all
observed worst-task excess is in principle recoverable, and separates the
cohorts on conditioning instead. Both halves of that invite a test. The
first invites an encoder that recovers the slack, and we built one: a
rate-distortion encoder using incoherent orthogonal mixing and scalar
quantization, regularised with a single Tikhonov ridge term
(\S\ref{sec:achv}, implementation in Appendix~\ref{app:manifest}). On the
shared-initialisation cohorts our earlier draft measured, it wins
comfortably against every baseline. The second invites the question of
whether a four-order-of-magnitude difference in conditioning is visible to
anything a practitioner does.

This section reports what happened when we tested that result properly.
Every decision rule quoted below was committed to version control before
the corresponding compute ran, and the commit ordering is the audit trail.
The pre-registrations are reproduced verbatim in
Appendix~\ref{app:prereg}.

\subsection{The encoder's advantage does not survive}
\label{subsec:q1}

\paragraph{Pre-registered question.} Does the encoder's win over the
best baseline survive the move from the shared-initialisation cohort to
independently initialised ones? The rule, fixed in advance: for each base
model, compare the encoder against the champion baseline selected on the
mean over the three independent cohorts, with a tie threshold of $0.005$
nats and a one-directional noise gate at $2\times\mathrm{SE}$ that may
only downgrade a win or loss to a tie, never promote.

\begin{table}[t]
\centering
\small
\caption{Pre-registered replication on independently initialised cohorts.
Worst-task NLL excess in nats, mean over three cohorts. Positive $d$
favours the encoder. The gate is computed on the paired difference
$d$, whose standard deviation across cohorts is $0.0003$ to $0.0020$;
individual cells vary more, from $0.0001$ to $0.0034$ across all seven
methods, because the encoder and the champion move together from cohort to
cohort and the difference is quieter than either. The gate downgraded
nothing.}
\label{tab:q1}
\begin{tabular}{lccccc}
\toprule
base & rd-encoder ridge & champion baseline & $d$ & $2\times\mathrm{SE}$ & result \\
\midrule
Llama-3.1-8B & 0.0717 & 0.1052 \; (TVQ$_2$) & $+0.0335$ & 0.0010 & wins \\
Mistral-7B   & 0.0597 & 0.0543 \; (TIES)    & $-0.0055$ & 0.0024 & loses \\
Qwen2.5-7B   & 0.0143 & 0.0141 \; (TIES)    & $-0.0002$ & 0.0004 & ties \\
Yi-1.5-9B    & 0.0637 & 0.0487 \; (TIES)    & $-0.0150$ & 0.0011 & loses \\
\bottomrule
\end{tabular}
\end{table}

\paragraph{Outcome: it does not survive.} Table~\ref{tab:q1} gives one
win, one tie and two losses. Both pre-specified robustness lines agree.
Re-selecting the champion inside each cohort, the variant deliberately
biased against our method, returns the same one/one/two. Applying the
original single-cohort rule three times returns a weakened verdict on
\texttt{indep1} but a failure on \texttt{indep2} and \texttt{indep3}, so
the earlier optimistic reading was the best of three draws.

\paragraph{One arm of that table was measured on a different inference
path, and matching it moves a verdict.} We found this while checking an
unrelated cell and report it because it bears directly on Table~\ref{tab:q1}.
Of the seven methods in that matrix, the ridge-regularised encoder is the
only one evaluated through vanilla \textsc{transformers}; the six others run
a fused fast path. That is not a slip: at $\deff$ realisation the merged
adapter has rank $64$ while the baselines have rank $16$, and the fast path
does not support mixed-rank adapters, so the encoder had to take the other
route. But the two paths are not numerically identical. On the same
adapters and the same evaluation draw they differ by up to $0.0105$ nats,
which is twice our tie threshold and comparable to three of the four margins
in the table.

The rank-$16$ realisation of the same encoder runs the fast path, so it is
matched to the baselines on both counts and needs no new measurement:

\begin{center}
\small
\begin{tabular}{lccc}
\toprule
base & $d$ as published & $d$ matched & verdict \\
\midrule
Llama-3.1-8B & $+0.0335$ & $+0.0257$ & wins, both \\
Mistral-7B   & $-0.0055$ & $-0.0022$ & \textbf{loses $\to$ ties} \\
Qwen2.5-7B   & $-0.0002$ & $-0.0004$ & ties, both \\
Yi-1.5-9B    & $-0.0150$ & $-0.0154$ & loses, both \\
\bottomrule
\end{tabular}
\end{center}

Matched, the count is one win, two ties and one loss rather than one win,
one tie and two losses. Mistral was reported as a clear loss on a margin
that cleared the threshold by $0.0005$ nats. The conclusion of this section
does not change, and neither does the direction of any other base, but the
matched arm is the one that should be read.

\paragraph{Nor does the remaining win survive dropping the untrained
adapter.} \S\ref{subsec:scope} reports that one of our four adapters never
learned its task, and Appendix~\ref{app:t3} shows every geometric quantity
is unchanged when it is removed. The merge results deserve the same
treatment, so we re-ran the matrix at $T = 3$ on all three independent
cohorts:

\begin{center}
\small
\begin{tabular}{lcc}
\toprule
base & $d$ at $T = 4$ & $d$ at $T = 3$ \\
\midrule
Llama-3.1-8B & $+0.0257$ \; wins & $+0.0033$ \; \textbf{ties} \\
Mistral-7B   & $-0.0022$ \; ties & $-0.0012$ \; ties \\
Qwen2.5-7B   & $-0.0004$ \; ties & $+0.0009$ \; ties \\
Yi-1.5-9B    & $-0.0154$ \; loses & $-0.0053$ \; loses \\
\bottomrule
\end{tabular}
\end{center}

One win, two ties and one loss becomes \emph{zero} wins, three ties and one
loss. The Llama win, the only clear win this method has left anywhere in
the paper, falls from $0.0257$ to $0.0033$ nats and stops clearing the
gate. The rank-$64$ realisation disagrees, keeping Llama at $+0.0081$ while
moving Mistral from a loss to a tie, and two realisations of one method
disagreeing about the same base is itself a measure of how thin that win
is. The champion on Mistral changes from TIES to TVQ at $T = 3$, which we
report rather than absorb.

The $T = 4$ matrix was registered as primary and stays primary; we are not
replacing it with the arm that happens to be less flattering, any more than
we would the reverse. But levels are not comparable across $T$, since a
maximum over three tasks is not a maximum over four, so what transfers is
the verdict, and the verdict is that three separate controls have each
removed a piece of this method's advantage: independent initialisation,
matching the rank and inference path, and now removing the adapter that
never trained.

We are careful about how strongly to put this, because the margins do
not all support the same strength of claim. The encoder wins clearly on
Llama ($+0.0335$, well outside the gate) and loses clearly on Yi
($-0.0150$). The other two are much weaker: Mistral at $-0.0055$ clears
the $0.005$ tie threshold by $0.0005$, and Qwen at $-0.0002$ is a tie by
any reading. So the accurate summary is \textbf{one clear win, one clear
loss, and two bases where the encoder and a 2023 baseline are
indistinguishable} on a metric whose own resolution we estimate at about
$0.001$ nats (\S\ref{sec:discussion}). What does not survive is the
earlier claim that the encoder is the better method: on properly
initialised cohorts it is not distinguishable from TIES on half the
bases and worse on a quarter. The Llama win is not a rank artifact:
rank-matching the encoder to the baselines still leaves it at $0.0795$
against the champion's $0.1052$.

We think the honest reading is that the encoder's apparent advantage was
substantially a property of the degenerate cohort it was developed on. On
that cohort $\bar H$ is ill-conditioned by four orders of magnitude
(\S\ref{subsec:conditioning}), which is precisely the situation the
encoder's ridge term was introduced to handle, so it was tuned on the one
regime where its distinguishing component has work to do.
\S\ref{subsec:ridge} makes that reading concrete: the ridge gain on the
shared cohort is $0.19$ to $9.73$ nats and on the independent cohort
$0.00$ to $0.10$.

\subsection{Does the conditioning matter? Only for methods that solve}
\label{subsec:ridge}

\S\ref{subsec:conditioning} found a four-order-of-magnitude difference in
$\kappa(\bar H)$ between the two initialisation regimes. A difference that
size which nothing observable depends on would be a curiosity, so we tested
it, pre-registering the prediction before running the cells.

\paragraph{First, the null, and what it can actually detect.} We compare
every method on matched cohorts, same method, same base, shared against
independent initialisation, with three cohorts per arm so the
pre-registered gate can be applied: a cell counts as an effect only if the
mean difference exceeds $\max(0.005, 2\,\mathrm{SE})$. Across $20$
base-by-method cells, $15$ show no detectable effect, $5$ favour
independent initialisation, and none favours shared.

The number that matters alongside that is the smallest effect the design
could have found, because a null without one is not evidence of absence.
The gate sits at $0.0052$ nats in the median cell and $0.0109$ in the
worst, so the honest statement is \emph{no effect detectable at about
$0.01$ nats}, not \emph{no effect}. It also settles a tension in an
earlier version of this work, which quoted a mean difference of $0.0014$
nats while estimating the metric's own resolution at about $0.001$: that
difference is below the median gate, so it is a quantity this design
cannot resolve in either direction.

The five cells that do clear the gate are worth naming rather than
averaging away, because four of them are not spread at random. Three are
the subspace-alignment method, on Mistral, Qwen and Yi, at $0.0096$ to
$0.0129$ nats; the others are TIES on Llama at $0.0151$ and TVQ on Yi at
$0.0087$. Task arithmetic and DARE, the two methods closest to a plain
mean, show nothing anywhere. That ordering is the one the mechanism below
predicts: a method that explicitly operates on the subspace geometry is
the heuristic most exposed to a change in that geometry, and it gains from
the cohort whose subspaces are not near-collinear. The effect is small,
consistent in direction, and we do not build anything on it.

\paragraph{Then the mechanism, measured.} The explanation for the other
$15$ cells is that those methods never go near the ill-conditioned
direction. That was an assertion in an earlier version and is now a
measurement, and it needs no forward passes: every merge here is weight
arithmetic on the adapter factors. For each merged solution we report its
norm relative to the mean task delta, and the share of its energy lying in
the bottom quartile of $\bar H$'s eigen-directions, in the projected
coordinates of \S\ref{subsec:conditioning}.

\begin{table}[t]
\centering
\small
\caption{What each solution does with the ill-conditioned directions.
Amplification is $\|\Delta_{\mathrm{merged}}\|/\|\bar\Delta\|$; tail
fraction is the share of squared norm in the bottom quartile of $\bar H$'s
spectrum. Ranges are over the four base models. On the independent arm a
tail fraction near $0.21$ is what a generic direction gives when a quarter
of the dimensions holds a quarter of the energy, so the signature is the
$0.73$ to $0.80$ on the shared arm, not the baseline it is measured
against.}
\label{tab:mechanism}
\begin{tabular}{lcccc}
\toprule
 & \multicolumn{2}{c}{shared} & \multicolumn{2}{c}{independent} \\
solution & amplification & tail & amplification & tail \\
\midrule
mean task delta            & 1.00 & 0.000 & 1.00 & 0.21 \\
task arithmetic            & 0.99 & 0.000 & 0.86--0.92 & 0.20 \\
DARE                       & 0.98--0.99 & 0.000 & 0.86--0.90 & 0.20 \\
TIES                       & 1.98--2.04 & 0.000 & 2.04--2.22 & 0.21 \\
TVQ ($b{=}2$)              & 1.95--2.16 & 0.000 & 1.68--1.95 & 0.20 \\
\addlinespace
RegMean, minimal ridge     & 35.4--51.4 & 0.73--0.79 & 3.41--3.67 & 0.29 \\
our encoder, $\lambda = 0$ & 36.7--56.7 & 0.78--0.80 & 3.42--3.67 & 0.29 \\
our encoder, $\lambda^\star$ & 2.20--2.56 & 0.000 & 2.84--3.05 & 0.28 \\
\addlinespace
exact interpolator $\tauH$ & 43.4--64.1 & 0.62--0.66 & 3.96--4.00 & 0.28 \\
\bottomrule
\end{tabular}
\end{table}

Table~\ref{tab:mechanism} gives it. On the ill-conditioned cohort every
heuristic sits at one to two times the mean task delta with a tail
fraction that rounds to zero at three decimals, while both unregularised
solvers sit at $35$ to $57$ times with roughly three quarters of their
energy in the bottom quartile of the spectrum. They are the only methods
that go there, which is why they are the only methods the conditioning
reaches. One incidental check falls out: RegMean at minimal ridge and our
encoder at $\lambda = 0$ agree to two decimals on the independent arm
across all four bases, so two differently derived solvers land on the same
solution once the problem is well posed.

\paragraph{So we sweep the ridge, on two solvers.} To expose the
conditioning one needs a method that actually solves the system. Ours
does, and it carries a Tikhonov ridge, so sweeping $\lambda$ over seven
values, $0$, $0.01$, $0.03$, $0.05$, $0.13$, $0.30$ and $1.00$, on both
cohorts with rank pinned to $16$ in every cell isolates the effect.

\begin{table}[t]
\centering
\small
\caption{Worst-task NLL excess for the unregularised solve
($\lambda = 0$), and the gain from the best ridge. Conditioning is
invisible to the heuristics above and severe here.}
\label{tab:ridge}
\begin{tabular}{lccccc}
\toprule
base & \multicolumn{3}{c}{$\lambda = 0$} &
\multicolumn{2}{c}{ridge gain $L(0) - L(\lambda^\star)$} \\
 & shared & independent & ratio & shared & independent \\
\midrule
Llama-3.1-8B & 0.3081 & 0.0713 & $4.3\times$  & 0.2216 & 0.0000 \\
Mistral-7B   & 9.7799 & 0.1504 & $65\times$   & 9.7301 & 0.0952 \\
Qwen2.5-7B   & 0.1973 & 0.0256 & $7.7\times$  & 0.1867 & 0.0114 \\
Yi-1.5-9B    & 0.2490 & 0.1209 & $2.1\times$  & 0.2211 & 0.0904 \\
\bottomrule
\end{tabular}
\end{table}

Table~\ref{tab:ridge} gives the result. Unregularised, the same merge is
$2.1$ to $65$ times worse on the ill-conditioned cohort; on Mistral it
collapses outright to $9.78$ nats against $0.15$. The ridge gain is larger
on the shared arm on all four bases, which was the pre-registered
prediction and held.

The other half of the prediction failed, and we report it as failed. We
also predicted that $\lambda^\star$ itself would be at least ten times
larger on the shared arm. It is not: $\lambda^\star$ is \emph{identical}
across arms on Mistral, Qwen and Yi ($0.13$, $0.13$, $0.30$), separating
only on Llama. How much regularisation is needed does not track
conditioning in any simple way; how much it buys does. By the
pre-registered rules this makes the test \textsc{partial}, and we claim
only the half that held.

\subsection{Both arms at three cohorts}
\label{subsec:shared-arm}

Table~\ref{tab:ridge} rests on one shared cohort against three independent
ones, so its noise gate borrowed its standard error entirely from the
independent side. A shared-versus-independent contrast measured that way is
three draws against one, and a reader is right to discount it. We registered
the completion in advance and ran the sweep on two further shared cohorts, so
both arms are $n = 3$ and the gate carries a term from each. This is a
replication of a result we had already published, so it could only cost us
something; the registration says so and fixes the consequence text for the
outcome where it fails.

It does not fail. With
$\mathrm{SE} = \sqrt{\mathrm{SE}_s^2 + \mathrm{SE}_i^2}$ and the gate at
$\max(0.005, 2\,\mathrm{SE})$, the ridge gain is larger on the shared arm on
\textbf{four of four} bases and the unregularised excess is worse there by at
least a factor of two on \textbf{four of four}: ratios of $4.42$, $76.2$,
$10.8$ and $2.82$, against gates of $0.024$, $4.26$, $0.128$ and $0.166$.
Mistral's shared arm averages $11.72$ nats at $\lambda = 0$ against $0.154$
independent.

One quantity is worth reporting because it could have embarrassed the earlier
table and does not. The published single-cohort value, \texttt{seed1}, sits
$0.23$, $0.53$, $0.75$ and $0.66$ standard deviations from its own arm's
three-cohort mean. We registered that a distance above $2$ on two or more
bases would oblige us to say the published value was unrepresentative of its
own arm. It is representative, and Table~\ref{tab:ridge}'s numbers stand as
single-cohort estimates of quantities the three-cohort means confirm.

\subsection{The same test on a solver we did not build}
\label{subsec:regmean}

Everything above concerns one solver, and we built it. A referee is
entitled to ask whether the effect is a property of the family or of our
implementation, and the question is sharp enough to be worth answering
with a pre-registered experiment rather than an argument.

RegMean~\citep{jin2023regmean}, in the data-free adapter-only form, is a
solver we did not design, whose ridge term is part of its own published
formulation rather than something we added. We swept it on the same grid
and the same two cohorts, registering two predictions in advance: that the
ridge gain would be larger on the shared arm, and that the minimally
regularised excess would be worse there by at least a factor of two, each
on at least three of four bases.

\begin{table}[t]
\centering
\small
\caption{RegMean on both cohorts. Its minimally regularised reference is
$\lambda = 10^{-6}$ rather than $0$: RegMean solves in the full input
dimension against a Gram of rank at most $Tr = 64$ against $d = 4096$, so
$\lambda = 0$ is singular rather than ill-conditioned and returns
null-space noise without raising an error. Our encoder solves inside
$\range(\bar H)$, so its $\lambda = 0$ is defined. The two columns are
therefore not the same object and each method is compared only against
itself across arms.}
\label{tab:regmean}
\begin{tabular}{llcccc}
\toprule
base & cohort & $\lambda = 10^{-6}$ & $\lambda^\star$ & gain & ratio \\
\midrule
Llama-3.1-8B & shared      & 1.9789 & 0.01 & 1.7428 & \multirow{2}{*}{$27.9\times$} \\
             & independent & 0.0708 & $10^{-6}$ & 0.0000 & \\
\addlinespace
Mistral-7B   & shared      & 10.6283 & 0.01 & 10.5314 & \multirow{2}{*}{$70.9\times$} \\
             & independent & 0.1499 & 0.01 & 0.0449 & \\
\addlinespace
Qwen2.5-7B   & shared      & 0.2376 & 0.01 & 0.2255 & \multirow{2}{*}{$9.3\times$} \\
             & independent & 0.0256 & 0.01 & 0.0133 & \\
\addlinespace
Yi-1.5-9B    & shared      & 0.5700 & 0.01 & 0.5477 & \multirow{2}{*}{$4.7\times$} \\
             & independent & 0.1203 & 0.01 & 0.0913 & \\
\bottomrule
\end{tabular}
\end{table}

Both predictions hold on four of four bases. The effect is not a property
of our encoder.

Three details in Table~\ref{tab:regmean} are worth reporting rather than
smoothing. $\lambda^\star$ is $0.01$ on seven of the eight arms, so on
RegMean too the amount of regularisation needed does not track the
conditioning even though the amount it buys does; that is the same split
our own sweep found, now reproduced on a method we did not write. Llama's
independent arm selects the minimal $\lambda$ and gains exactly zero,
which is the cleanest statement of the claim available: on a
well-conditioned cohort there is nothing to repair. And from
$\lambda = 0.13$ upward the two arms agree to three decimals, because
heavy regularisation washes the difference out entirely.

\subsection{The same test without rank truncation}
\label{subsec:untruncated}

\S\ref{subsec:conditioning} notes that our pipeline truncates every merge
back to rank $r$, so the sweep above tests conditioning through a
rank-constrained solver while the interpolator whose norm we report has
rank up to $Tr$. Repeating the sweep with truncation disabled removes that
mismatch, and the truncated sweep is then the comparison arm rather than
being re-run.

\begin{center}
\small
\begin{tabular}{lcccc}
\toprule
 & \multicolumn{2}{c}{$\lambda = 0$ excess} & \multicolumn{2}{c}{shared / independent} \\
base & shared & independent & untruncated & truncated \\
\midrule
Llama-3.1-8B & 7.4266  & 0.0635 & $116.9\times$ & $4.3\times$ \\
Mistral-7B   & 11.6283 & 0.1582 & $73.5\times$  & $65\times$ \\
Qwen2.5-7B   & 0.8993  & 0.0256 & $35.2\times$  & $7.7\times$ \\
Yi-1.5-9B    & 3.1294  & 0.1227 & $25.5\times$  & $2.1\times$ \\
\bottomrule
\end{tabular}
\end{center}

The link survives on all four bases, and the effect is \emph{larger}
without truncation, by factors of $25$ to $117$ against $2.1$ to $65$.
Truncation was not manufacturing the result; it was hiding part of it,
which is what one should expect when all but the leading $16$ directions
of a rank-$64$ interpolator are discarded and the ill-conditioned subspace
goes with them. The numbers we report as primary are therefore the
conservative ones.

\paragraph{What this licenses, and what it does not.} It licenses a
conditional recommendation: if you merge with a method that solves a
linear system, check your cohort's conditioning and regularise, because
unregularised it can fail by a factor of $65$ at rank $16$ and $117$
without truncation, on two solvers, one of which we did not build. It does
\emph{not} license the general advice to initialise independently. On the
methods most practitioners actually use the null above stands, with the
single qualification that the subspace-alignment method gains a little
from independent initialisation on three of four bases, and we do not
dress that up as anything more.

\subsection{A claim we withdraw}
\label{subsec:withdrawn}

The result in \S\ref{subsec:q1} raises an obvious and much more
interesting possibility: if method rankings shift when initialisation
changes, then published merging benchmarks, which as far as we can tell
mostly use shared-seed cohorts, might be measuring initialisation
geometry rather than method quality. We described this on first sight as
the strongest result available to us.

We pre-registered a control to test it, and the control did not support
the claim. The control asks whether the method ranking is stable
\emph{within} the independent regime, across the three cohorts that
differ only in initialisation draw. If rankings are unstable within a
single regime, then ranking changes across regimes cannot be attributed
to the regime. Two of four base models had a non-unanimous top-1 method
across the three cohorts, above the threshold fixed in advance, so by our
own rule the attribution fails and \textbf{the claim that published
merging benchmarks are confounded by shared-initialisation geometry is
withdrawn}.

We record one caveat, and we deliberately did not let it reverse the
verdict. The control as written counts a top-1 change as instability
regardless of margin, and both flips involve methods separated by far less
than the tie threshold: on Qwen the top two are $0.0140$ and $0.0141$, a
coin flip. A margin-aware version of the control would very likely give a
different answer. But that version was not what we pre-registered, and
rewriting a decision rule after seeing which way it fell is exactly the
practice the pre-registration exists to prevent. We stated the defect and
left the verdict, and registered a \emph{second} control rather than
editing the first.

\paragraph{The second control, and what it does and does not restore.} The
margin-aware rule was fixed in the same registration that governs
\S\ref{subsec:regmean}: a base counts as unstable only if its top-1 method
changes \emph{and} the margin between the top two exceeds the tie threshold
in at least one cohort where the change occurs. It needs no new cells. Run
as written it returns \textbf{zero of four} bases unstable, against two of
four for the margin-blind rule. Both flips are inside the threshold by a
wide margin: Mistral's are $0.0004$, $0.0003$ and $0.0011$ nats, Qwen's
$0.0001$, $0.0006$ and $0.0003$, against a threshold of $0.005$.

By the registered rule that reinstates the benchmark-confounding claim as
\textsc{provisional}. Three things bound what that means, all of them fixed
in advance. The margin-blind control was registered first and is not
deleted or replaced; it is reported here alongside. The margin-aware control
was registered after its data had been seen, which the registration states
about itself, so it is secondary evidence and we treat it as such. And the
parent registration attached a precondition to any claim about the field,
which is that the shared-initialisation pattern must be shown to occur
outside our own repository before we assert anything about published
benchmarks. That precondition is now discharged by
\S\ref{subsec:prevalence}, where $37.8\%$ of $45$ public cohorts are in the
collapsed regime.

So the claim we make is this, and no more: method rankings on this metric
move between cohorts that differ only in initialisation draw, the movement
is within noise once margin is accounted for, and the collapsed regime that
would make such rankings hard to interpret is common in released adapters.
We have not measured a published merging benchmark, and we do not claim one
is wrong.

What survives independently of all of it is weaker and still worth saying:
single-seed method rankings on this metric are not reproducible, so merging
comparisons should report multiple initialisation draws.

\subsection{The collapse is visible in downstream accuracy}
\label{subsec:downstream-cond}

Everything above is measured in worst-task NLL excess, and
\S\ref{sec:discussion} reports that this metric does not predict downstream
accuracy across merge methods, and on HumanEval predicts it backwards. A
reader is entitled to extend that objection to the whole paper: if the metric
does not track model quality, why should any of these numbers matter?

The objection is about differences of hundredths of a nat between methods. The
conditioning effect is not that. At $\lambda = 0$ the shared-versus-independent
gap runs from $0.13$ to $11.7$ nats, one to three orders of magnitude above
both the metric's resolution and the differences the correlation was computed
over. An effect that size either shows up in accuracy or the metric is
measuring something other than model quality. We registered the test without
knowing which, and fixed in advance that a null would move this limitation into
\S\ref{sec:intro} and delete the practice recommendations rather than defend
them.

\begin{table}[t]
\centering
\small
\caption{Downstream accuracy of the same merges, with the repaired scorers.
GSM8K exact match at $n = 500$, HumanEval pass@1 at $n = 164$, its whole test
split. Gate is $\max(5\text{ points}, 2\,\mathrm{SE})$ on the binomial. No cell
had a scorer discard: every generation was scored.}
\label{tab:downstream-cond}
\begin{tabular}{lcccccc}
\toprule
& \multicolumn{3}{c}{shared} & \multicolumn{3}{c}{independent} \\
base & $\lambda{=}0$ & $\lambda^\star$ & gain
     & $\lambda{=}0$ & $\lambda^\star$ & gain \\
\midrule
\multicolumn{7}{l}{\emph{GSM8K, exact match}} \\
Llama-3.1-8B & 0.318 & 0.654 & $+0.336$ & 0.708 & 0.712 & $+0.004$ \\
Mistral-7B   & \textbf{0.000} & 0.462 & $+0.462$ & 0.276 & 0.424 & $+0.148$ \\
Qwen2.5-7B   & 0.496 & 0.808 & $+0.312$ & 0.786 & 0.786 & $+0.000$ \\
Yi-1.5-9B    & 0.346 & 0.802 & $+0.456$ & 0.698 & 0.794 & $+0.096$ \\
\addlinespace
\multicolumn{7}{l}{\emph{HumanEval, pass@1}} \\
Llama-3.1-8B & 0.250 & 0.573 & $+0.323$ & 0.518 & 0.530 & $+0.012$ \\
Mistral-7B   & \textbf{0.000} & 0.244 & $+0.244$ & 0.067 & 0.189 & $+0.122$ \\
Qwen2.5-7B   & 0.530 & 0.744 & $+0.213$ & 0.726 & 0.738 & $+0.012$ \\
Yi-1.5-9B    & 0.024 & 0.091 & $+0.067$ & 0.049 & 0.067 & $+0.018$ \\
\bottomrule
\end{tabular}
\end{table}

\paragraph{It shows up, on both benchmarks.} The ridge improves the shared arm
past the gate on \textbf{four of four} bases on both benchmarks. The
difference-in-differences, which is the claim that the ridge helps
\emph{because of} conditioning rather than helping everything everywhere,
clears the gate on four of four on GSM8K and three of four on HumanEval; the
exception is Yi, where both arms sit near the floor of a benchmark it barely
solves at all ($0.024$ to $0.091$). By the registered rule, requiring either
prediction on at least three of four bases on at least one benchmark, the
outcome is \textsc{confirmed}.

\paragraph{The number worth stating plainly.} On Mistral, the shared-cohort
unregularised merge scores \textbf{zero} on both benchmarks. It is not
producing malformed output that a scorer discards: all $500$ GSM8K generations
and all $164$ HumanEval generations are non-empty, and none contains an
extractable answer or passes a test. The model generates fluent text and has
lost the capability entirely. One Tikhonov term on the same adapters recovers
$0.462$ and $0.244$. That is what four orders of magnitude of conditioning buys
in a unit a practitioner cares about.

\paragraph{What this does not license.} It does not show that worst-task NLL
excess predicts accuracy in general, and we do not claim it. The correlation
that fails in \S\ref{sec:discussion} is across merge \emph{methods} at fixed
cohort, and nothing here tests that. What it shows is that the specific effect
this paper is about is not an artifact of the metric, which is the objection it
was registered to answer.

\subsection{The predicted rate exponent is unresolved}
\label{subsec:w3}

Theorem~\ref{thm:general} predicts distortion decaying as
$2^{-2R/\deff}$, which on a log-rate axis is a slope in a narrow band. We
pre-registered the band as $[-2.4, -1.6]$, to hold on at least three of
four base models, and swept quantizer width $b \in \{1,2,3,4,8,16\}$ with
the fit taken over $b \in \{2,3,4,8\}$, $b=1$ being clipping-dominated.

An earlier version of this work reported the prediction falsified, on the
strength of measured slopes of $-0.10$ and $-0.21$. \textbf{Those two
numbers are withdrawn.} Each finite-rate curve was compared against a
$b = \infty$ reference computed at a \emph{different} ridge value and
aggregated over a different number of cohorts, so the quantity being
differenced was not the same method. The symptom was visible in the
output and we did not act on it at the time: on two bases the excess at
finite rate came out \emph{below} its own asymptote, which the model does
not permit. That is a measurement fault, not a property of the data.

\begin{table}[t]
\centering
\small
\caption{Rate sweep against matched asymptotes: same ridge $\lambda$, same
realisation, same cohort for the curve and its $b = \infty$ reference, with
the equality asserted from the configuration files before any fit is
attempted. Worst-task NLL excess in nats, at four decimals because at three
the reader cannot check the two facts that matter. \textbf{Ten of the
twenty-four measured values fall below their own asymptote}, which the model
forbids; all but one are smaller than $10^{-4}$ nats and are consistent with
numerical noise, the exception being Llama at $b=3$ at $-0.0020$. The fit is
ordinary least squares of $\log_2(\text{excess} - \text{asymptote})$ on $b$
over $b \in \{2,3,4,8\}$, keeping residuals above $10^{-9}$; the ``fit''
column names the points that survived that filter, and both surviving fits
rest partly on residuals below this metric's own resolution of about
$0.001$ nats (\S\ref{sec:discussion}).}
\label{tab:w3}
\begin{tabular}{lccccccccccc}
\toprule
base & $b{=}1$ & $b{=}2$ & $b{=}3$ & $b{=}4$ & $b{=}8$ & $b{=}16$ &
$b{=}\infty$ & fit & slope & $R^2$ \\
\midrule
Llama-3.1-8B & 0.6329 & 0.0967 & 0.0803 & 0.0845 & 0.0823 & 0.0823 & 0.0823 & 2,4,8 & $-2.10$ & 0.9843 \\
Mistral-7B & 0.2106 & 0.0497 & 0.0480 & 0.0471 & 0.0474 & 0.0473 & 0.0474 & 2,3 & --- & --- \\
Qwen2.5-7B & 0.0746 & 0.0110 & 0.0102 & 0.0101 & 0.0101 & 0.0100 & 0.0102 & 2,3 & --- & --- \\
Yi-1.5-9B & 0.2099 & 0.0413 & 0.0361 & 0.0357 & 0.0356 & 0.0355 & 0.0356 & 2,3,4 & $-3.50$ & 0.9996 \\
\bottomrule
\end{tabular}
\end{table}

\paragraph{Refitted, the verdict is unresolved.} Table~\ref{tab:w3} gives
the corrected sweep. Llama fits at $-2.10$ with $R^2 = 0.98$ against a
predicted $-2$, and Yi at $-3.50$, steeper than predicted rather than
flatter. Mistral and Qwen cannot be fitted, because only two of their four
in-window points have a residual above the filter at all.

\paragraph{A second fault in this analysis, found in preparing this
version.} The first fault was the mismatched reference described above. The
second is that neither surviving fit is supported by the data at the
precision the metric provides, and the published three-decimal table
concealed it.

The fits are on the residual above the asymptote, and those residuals are
mostly tiny. Llama's three fitted points have residuals $0.0144$, $0.0023$
and $0.0000029$ nats. The last is the $b = 8$ point, and it is what makes
the fit possible at all: without it Llama has two usable points and is
unfittable like Mistral and Qwen. It sits about $350$ times below this
metric's own stated resolution of $0.001$ nats
(\S\ref{sec:discussion}), so on a $\log_2$ axis it carries enormous
leverage while being, by our own account of the metric,
indistinguishable from zero. Yi is worse in one respect: two of its three
fitted residuals, $0.00046$ and $0.000044$, are below the resolution.

We also withdraw the account this section previously gave of why Mistral
and Qwen are unfittable. It said each ``reaches its asymptotic value
inside the fit window'', which is benign and wrong. What actually happens
is that their $b=4$ and $b=8$ points fall \emph{below} their asymptotes and
are removed by the residual filter. That is the same anomaly this section
opened by describing, and the caption of the published table asserted it no
longer occurred. It occurs on all four bases, ten times in twenty-four
measurements, and Table~\ref{tab:w3} now reports it.

The honest reading is that this sweep does not have the dynamic range to
measure a rate exponent. Only $b \in \{1, 2\}$ is comfortably above the
noise floor on any base, and $b = 1$ is excluded by the registered rule as
clipping-dominated. We are therefore not entitled to the claim that the
exponent holds on Llama, and we withdraw it; it was the only positive
evidence for the rate term anywhere in this paper.

The rule fixed in advance requires the slope to lie in $[-2.4, -1.6]$ on
at least three of four bases to void the falsification, and $|{\rm slope}|
< 0.5$ on at least three of four for it to stand. We obtain one and zero
respectively, so by the registered rule the outcome is
\textsc{unresolved}. We apply the rule as written, to the fits as computed,
rather than adjusting the residual filter now that we can see which points
it admits: moving a threshold after seeing its consequences is the practice
the pre-registration exists to prevent, and it would in this case move the
verdict in our own favour. Note however that the reading above makes the
outcome unresolved for a more basic reason than the rule contemplates. The
rule presumes fits worth judging, and on this grid there are none. The
registered window $\{2,3,4,8\}$ is too coarse for bases that reach their
noise floor by three bits; a finer low-rate grid would settle it, and we
have deliberately not chosen one after seeing which grid would produce a
fit.

What survives is an operational regularity rather than a test of the
theory, and it is the one thing here the precision does support. With the
ridge on, three bits buys essentially the whole effect: from $b = 3$ upward
every base is within $0.005$ nats of its own asymptote, and sixteen bits
adds nothing measurable over four. Only $b=1$, where clipping dominates, is
far off. A practitioner quantizing a merged adapter can stop at three or
four bits, and that recommendation does not depend on any exponent.

Three caveats. This sweep is a single initialisation draw on the
shared-initialisation cohort, so it says nothing about whether the
finite-rate behaviour differs by regime; repeating it on an independent
cohort would be a genuinely new measurement and needs its own registered
rule, since it asks a different question from the one registered here.
The faults corrected above were both ours and both in the analysis rather
than the cells: the first in the reference the analyser selected, the
second in the precision at which we reported and read the result. And we
note that the second fault was invisible at three decimals and would have
stayed invisible: a reader could not have derived it from the published
table, which is the reason we now print the fitted points and the residual
scale rather than the rounded values alone.

\subsection{DARE: a negative result with a mechanism, and its limits}
\label{subsec:dare}

DARE~\citep{yu2024dare} drops a random fraction of task-vector entries
and rescales the survivors by $1/(1-p)$. Composed with task arithmetic it
does essentially nothing in our setting: across a density sweep from
$0.05$ to $0.5$ on four base models, it beats plain task arithmetic in
three of twenty cells and all three are within $0.0005$ nats, a fifth of
our tie threshold.

Unlike most negative results this one has a mechanism that predicts its
own shape. The rescaling makes the masked task vector an unbiased
estimator of the unmasked one, $\E[D_{\mathrm{DARE}}] = D_{\mathrm{TA}}$.
We state the consequence carefully, because the obvious version of the
argument is not available here: negative log-likelihood of an LLM is
\emph{not} convex in the weights, so Jensen's inequality does not apply
globally and we do not claim
$\E[L(D_{\mathrm{DARE}})] \geq L(D_{\mathrm{TA}})$ as a theorem. What
does hold is the local statement. Expanding $L$ to second order about
$D_{\mathrm{TA}}$, the first-order term vanishes by unbiasedness and the
penalty is $\tfrac12\tr(H\Sigma) + o(\|\Sigma\|)$ with
$\Sigma \propto (1-p)/p$ and $H$ the local Hessian. Where $H \succeq 0$,
which is the generic situation near a trained solution, this is
non-negative and grows as density falls. That is what we observe,
on three of four base models. We verified our implementation against the
reference one and confirmed the operator reproduces the analytic signature
$\sqrt{(1-p)/p}$ to three decimals, so this is a property of the method
and not of our code.

\paragraph{The negative result does not generalise, and we tested that.}
The Jensen argument applies only to the task-arithmetic composition, since
TIES is biased and the mask changes which coordinates survive trimming. We
pre-registered the DARE-TIES composition test before implementing it. The
outcome is genuinely mixed: DARE composed with TIES \emph{helps}
decisively on Llama ($+0.0462$ nats, against a gate of $0.0036$) and hurts
decisively on the other three. By the rule fixed in advance this is a
mixed verdict and neither direction is claimed.

So the correct scope for the negative result is DARE composed with task
arithmetic, not DARE. We flag that Llama is now the odd model out in two
separate pre-registered tests, here and in Table~\ref{tab:q1}, without
proposing a mechanism for it; a base-level moderator that neither
pre-registration anticipated deserves a designed experiment rather than
another post-hoc reading.
